% Created 2026-08-26 Wed 17:31
% Intended LaTeX compiler: pdflatex
\documentclass[11pt]{article}
\usepackage[utf8]{inputenc}
\usepackage[T1]{fontenc}
\usepackage{graphicx}
\usepackage{longtable}
\usepackage{wrapfig}
\usepackage{rotating}
\usepackage[normalem]{ulem}
\usepackage{amsmath}
\usepackage{amssymb}
\usepackage{capt-of}
\usepackage{hyperref}
\author{fcalado}
\date{\today}
\title{}
\hypersetup{
 pdfauthor={fcalado},
 pdftitle={},
 pdfkeywords={},
 pdfsubject={},
 pdfcreator={Emacs 29.3 (Org mode 9.6.15)}, 
 pdflang={English}}
\begin{document}

\tableofcontents

\section{Class notes}
\label{sec:orgaa5d2ae}

\subsection{Week 1}
\label{sec:org5e42edd}

\begin{enumerate}
\item In your understanding, why do automated processes discriminate
against certain groups, what factors specifically causes them to
discriminate? Try to guess the factors.
\begin{itemize}
\item it could be the model is not a equitable choice.
\item it could be the dataset — there’s more training samples (existing
data) of majority groups having jobs, so the model will
perpetuate these trends
\item scaling up creating a cycle issue of bias
\item aggregate analyses can make it difficult to see specific
instances of discrimination
\end{itemize}

\item Do you perceive any potential oversights or flaws in the authors’
arguments? Is there anything that should be interrogated or
explored in more depth?
\begin{itemize}
\item leaves a considerable aspect unaddressed: does not consider
applicants with Visas, who are automatically turned down.
\item could also expand into other aspects of discrimination like
gender, class, other aspects of background.
\item what about the other vendors?
\item could introduce some controls to see if the results change?
\item may not be considering how multiple apps / person might create
some bad fits for job apps.
\end{itemize}
\end{enumerate}

Python challenge:

Write some python code that does something while containing the
following components listed below.
\begin{itemize}
\item a loop
\item a list
\item a function
\end{itemize}

Solution:
\begin{verbatim}
import random

def greeting(names):
    greetings = [
        "how are you doing today?",
        "tudo bem?",
        "what's shakin'?",
        "oi!",
        "como te sientes?"
    ]
    for i in names:
        print(f"Hello {i}, {random.choice(greetings)}")


greeting(['Kirsten', 'Filipa', 'Tony', 'Clare'])
\end{verbatim}

Output:
\begin{verbatim}
Hello Kirsten, como te sientes?
Hello Filipa, oi!
Hello Tony, what's shakin'?
Hello Clare, how are you doing today?
\end{verbatim}
\end{document}
